This thesis set out to compare two approaches to AI-agent-assisted software development --- an unstructured, single-agent approach and a structured, multi-agent approach --- by building the same application under each, and to derive from that comparison a set of best practices for human-agent collaboration in software engineering. Across both parts, the comparison showed that the structure imposed on a collaboration, not the underlying model, was the more decisive factor in the quality and consistency of the resulting software: the same model produced markedly different outcomes depending on whether it operated under explicit architectural rules, mandatory review gates, and formalized testing discipline, or without them.

Between the two, the structured approach is, in my view, clearly the stronger one. It caught problems earlier, when they were still cheap to fix, rather than after they had already spread across the codebase; it produced code and tests that held up more consistently against the same evaluation criteria; and it gave human oversight a fixed, comparable structure to act on rather than an open-ended judgment call made after the fact. The unstructured approach was not without merit --- it was faster to start and required less upfront investment --- but the advantages of the structured approach outweighed that cost throughout this project.

A few things, in hindsight, could have been approached differently, and are worth naming as directions for future work rather than as part of this thesis's own scope. Documentation consistency, addressed here through a hand-written specification paired with an automated validation test, could instead be achieved through tooling that observes an application's real traffic --- including the traffic a test suite already produces --- and generates the specification directly from observed behavior, removing the need to maintain either annotations or a separate check. The comparison itself could also be strengthened by a study design that varies the order in which the two approaches are run, or that holds the development environment constant across both, to separate the effect of structure more cleanly from the effects of learning and tooling that this thesis's single-pass design could not fully isolate. Finally, this project was one developer working on one codebase, continuously, from start to finish; whether the same structure holds up with multiple developers, multiple agents working concurrently, or a larger, longer-lived codebase with more competing conventions to enforce, remains untested and is a natural direction for future work to take up.
